% A subseção "Fundamentos do Paralelismo" começa aqui.
\subsecao{Fundamentals of Parallelism}

Parallelism in computing refers to the simultaneous execution of multiple tasks or parts of the same task, with the primary goal of reducing a program's total execution time. It is crucial to distinguish parallelism from concurrency. Concurrency occurs when multiple tasks are in progress at a given time, but may not be executing simultaneously. A single-core system can exhibit concurrency by rapidly switching between the execution of different tasks, simulating parallelism through preemption. True parallelism, on the other hand, requires hardware with multiple processing units \cite{Butenhof1997}. This distinction is fundamental to the present work, which exploits true hardware parallelism through multiple CPU cores. % review: comment 007 (issue #102)

\subsubsecao{Parallel Computing Architectures}

Parallel computer architectures are classified according to various criteria. A prominent taxonomy, proposed by Flynn \cite{Flynn1966}, categorizes these architectures based on instruction and data streams:
\begin{itemize}
    \item \textbf{SISD (Single Instruction, Single Data)}: Represents a conventional single-processor architecture, where a single instruction operates on a single data stream;
    \item \textbf{SIMD (Single Instruction, Multiple Data)}: Involves multiple processing units that execute the same instruction concurrently on different data elements;
    \item \textbf{MISD (Multiple Instruction, Single Data)}: Characterizes systems where multiple processing units apply distinct instructions on the same data stream;
    \item \textbf{MIMD (Multiple Instruction, Multiple Data)}:  Involves multiple processing units that execute distinct instructions on distinct data streams independently. This paradigm is directly related to multi-core architectures.
\end{itemize}

Within the MIMD category, systems are further distinguished based on their memory organization \cite{Pacheco2011}:
\begin{itemize}
    \item \textbf{Shared-Memory Systems}: In these systems, multiple processors share a physical memory address space. Inter-processor communication is typically achieved by reading from and writing to shared variables in memory. Multi-core processors are the most common and widely available example of this architecture, where all cores on a chip share access to main memory;
    \item \textbf{Distributed-Memory Systems}: Each processor in these systems has its own locally accessible memory, and communication between processors is usually done by exchanging messages across an interconnection network. Computer clusters are examples of this architectural model.
\end{itemize}

The present work focuses on the acceleration of algorithms in environments with multi-core processors, which fall into the shared memory MIMD category. A multi-core processor (MCP) is a single chip that contains two or more central processing units (CPUs) called ‘cores’. Each of these cores is capable of reading and executing program instructions independently, essentially functioning as a complete processor.

\subsubsecao{Processes vs Threads}

Within shared-memory systems, such as multi-core processors, software-level parallelism is commonly achieved through the execution of multiple threads within a single process.

A process is an instance of an executing program, an entity possessing its own dedicated virtual address space, file descriptors, environment variables, and other Operating System (OS) resources, all isolated from other processes \cite{Pacheco2011}.

A thread, on the other hand, is a lighter unit of execution started by a process. The same process can control, concurrently or in parallel, multiple threads. A key characteristic is that all threads within a process share its virtual address space, which encompasses executable code, global data, and the heap segment \cite{Pacheco2011}. However, each thread maintains its private resources, which are essential for its independent execution context:
\begin{itemize}
    \item A unique identifier (thread ID);
    \item A set of CPU registers, including the program counter (PC);
    \item A separate execution stack, used for local variables and function call control;
    \item Thread-specific state, such as signal masking and scheduling priority.
\end{itemize}

\subsubsecao{POSIX Thread Model (Pthreads)}

The POSIX Thread Model, commonly referred to as “Pthreads”, is a standard that defines an application programming interface (API) for thread creation and management, as specified by IEEE Std 1003.1c-1995. Pthreads provides a portable suite of data types, functions, and constants that allows developers to create multithreaded applications that are compilable and executable across diverse POSIX-compliant operating systems (e.g., Linux, macOS, and various Unix variants). The Pthreads API provides functionality for managing the thread lifecycle—including creation, termination, and joining—and, critically, incorporates synchronization primitives. These primitives, such as mutexes and condition variables, are essential for managing access to shared resources within shared-memory environments \cite{Butenhof1997} \cite{Pacheco2011}.

\subsubsecao{Cache Architecture and Memory-Bound Parallelism}

Modern shared-memory processors organize their caches in a hierarchy. Each core has private L1 and L2 caches—typically tens to hundreds of kilobytes—for low-latency access to recently used data. A shared last-level cache (LLC), often several megabytes, is accessible to all cores on the chip and serves as the final caching tier before main memory \cite{Pacheco2011}.

A workload is \textit{compute-bound} when its execution time is dominated by arithmetic operations, leaving memory bandwidth underutilized. It is \textit{memory-bandwidth-bound} (or simply \textit{memory-bound}) when the bottleneck is the rate at which data can be fetched from DRAM rather than computational throughput. Pattern-matching algorithms that operate on large automata—whose transition tables exceed the LLC capacity—fall into this regime: each state transition triggers a cache miss that must be resolved from DRAM, making throughput a function of available memory bandwidth rather than CPU frequency or core count \cite{Grama2003}.

Parallelism can intensify memory pressure. When multiple threads simultaneously traverse different regions of a large, read-only automaton, the combined working set exceeds the LLC, raising the aggregate cache-miss rate for all threads. Two additional hazards arise when threads share writable data. \textit{Cache coherence} protocols (such as MESI) ensure that modifications made by one core become visible to all others; they introduce inter-core invalidation traffic whenever a shared cache line is written \cite{Pacheco2011}. \textit{False sharing} occurs when two threads write to logically independent variables that happen to reside on the same cache line, causing coherence invalidations despite the absence of a true data dependency. False sharing can be avoided by aligning per-thread accumulators to cache-line boundaries.

Non-Uniform Memory Access (NUMA) architectures differentiate memory latency based on which socket owns the physical memory being accessed. The experimental platform used for the main measurements in this work is a single-socket Uniform Memory Access (UMA) machine; nevertheless, NUMA is an important consideration in multi-socket server deployments where memory-bound workloads exhibit significant latency asymmetry \cite{Grama2003}.

\subsubsecao{Hybrid P/E-Core Architectures}

Modern consumer and mobile processors have introduced a \textit{hybrid}
microarchitecture that integrates two distinct core types on a single die. A
typical Alder Lake mobile processor pairs high-performance cores (P-cores), often
with simultaneous multithreading, with energy-efficient cores (E-cores) that
prioritize throughput per watt.

P-cores feature wider out-of-order execution pipelines and higher operating
frequencies, delivering greater per-thread throughput at higher power
consumption. E-cores are physically smaller and optimized for
throughput-per-watt rather than single-thread latency; they operate at lower
frequencies. When equal-size work units are assigned across both core types, the
faster cores can finish early and wait at a synchronization barrier while slower
cores complete their share. This is the architectural reason for the
frequency-weighted and dynamic load-balancing strategies introduced later: they
are mechanisms for reducing stragglers on heterogeneous processors, not a change
to the Aho--Corasick matching semantics. The main experimental campaign of this
thesis is conducted on homogeneous cores; the hybrid-core discussion appears only
as a diagnostic comparison in the results chapter.
% Fim da subseção "Fundamentos do Paralelismo".

% A subseção "Parallel Performance Evaluation Metrics" começa aqui.
\subsecao{Parallel Performance Evaluation Metrics}

Performance evaluation is fundamental to understanding the effectiveness of a parallel program. The following metrics compare the execution time of the parallel program with its equivalent sequential counterpart \cite{Grama2003}. % review: comment 008 (issue #103) — relate the metrics to the experiments
The evaluation of parallel algorithms requires objective metrics that quantify both performance gains and resource utilization; accordingly, the metrics presented in this section provide the basis for evaluating the implementation proposed in this work and for interpreting the experimental results reported in Chapters 5 and 6.

\subsubsecao{Execution Time and Overhead}

The performance of a program is measured by its execution time. For parallel programs, the following terms are defined:
\begin{itemize}
    \item \textbf{Ts}: The execution time of a baseline sequential algorithm executed on a single processor unit.
    \item \textbf{Tp}: The execution time of the parallel program utilizing p processors.
\end{itemize}

Ideally, \texttt{Tp} is expected to be substantially less than \texttt{Ts}. However, parallel execution introduces \textit{overhead}. This encompasses the time spent by processors on tasks that do not directly contribute to problem resolution, including communication, synchronization, load balance, and other computations inherent to the parallel implementation \cite{Grama2003}.

\subsubsecao{Speedup}

Speedup (\texttt{S}) measures the performance gain achieved through parallel execution relative to sequential execution, as shown in Eq.~\ref{eq:speedup}. It is defined as the ratio of sequential execution time to parallel execution time:
\begin{equation}
    S = \frac{T_s}{T_p}
    \label{eq:speedup}
\end{equation}
An ideal, or linear, speedup would be \texttt{S = p}. Amdahl's Law~\cite{Amdahl1967} establishes a theoretical upper bound on the speedup achievable by parallelizing a program with a fixed problem size. If \texttt{s} represents the fraction of a program's sequential execution time that is inherently serial (non-parallelizable), and \texttt{f} is the parallelizable fraction (\texttt{s + f = 1}), the maximum speedup on \texttt{p} processors is given by Eq.~\ref{eq:amdahl}.
\begin{equation}
    S \leq \frac{1}{s + f/p}
    \label{eq:amdahl}
\end{equation}
In the limit, as the number of processors (\texttt{p}) tends to infinity, the maximum speedup is bounded by $\frac{1}{s}$. This implies that even a small sequential fraction (small \texttt{s}) can severely limit the total speedup, regardless of the number of processors available. Figure~\ref{fig:speedup-plot} illustrates the theoretical speedup according to Amdahl's Law compared to ideal speedup for a given sequential fraction. 

% review: comments 009 (issue #104) caption below figure, 010 (issue #105) no source line for author-made figures
\begin{figure}[H]
  \centering
  \includegraphics[width=0.7\linewidth]{speedup-plot.png}
  \caption{Ideal vs.\ Amdahl’s theoretical speedup ($s=0.1$).}
  \label{fig:speedup-plot}
\end{figure}

A complementary perspective is offered by the Gustafson-Barsis Law~\cite{Gustafson1988} in Eq.~\ref{eq:gustafson}, where $\alpha$ denotes the fraction of the \textit{parallel} run time spent in non-parallelizable (serial) work (and is therefore a distinct quantity from Amdahl's $s$, which is the serial fraction of the \textit{sequential} run). This law argues that, for problems whose size can be scaled with the available processors, an approximately linear speedup is attainable even in the presence of a serial fraction. Both Amdahl's and Gustafson's laws are important for understanding the limits and potential of parallelism in different scenarios \cite{Pacheco2011}.

\begin{equation}
    S_{\text{scaled}} = \alpha + (1-\alpha) \cdot p
    \label{eq:gustafson}
\end{equation}

\subsubsecao{Efficiency}

Efficiency (\texttt{E}) measures the effectiveness of processor utilization in a parallel program. It is defined as the ratio of speedup to the number of processors \texttt{p}, as in Eq.~\ref{eq:efficiency}:
\begin{equation}
    E = \frac{S}{p} = \frac{T_s}{p \cdot T_p}
    \label{eq:efficiency}
\end{equation}
Efficiency values range between 0 and 1 (inclusive), often expressed as a percentage (0\% to 100\%). An efficiency of 1 (or 100\%) signifies ideal, linear speedup, implying that all processors are contributing productively throughout the execution. In practice, efficiency typically diminishes as the processor count increases, primarily due to the escalating impact of parallel overhead \cite{Grama2003}.


\subsubsecao{Scalability}

In parallel computing, scalability denotes a system’s ability to sustain high performance (often via efficiency) as both the number of processors and/or the problem size grow \cite{Grama2003}.

\begin{itemize}
    \item \textbf{Strong Scalability}: Measures how execution time changes when the processor count $p$ increases for a fixed total problem size, denoted $W$ (for example, the total number of array elements to process). The goal is to keep efficiency nearly constant as $p$ increases. Amdahl’s Law directly applies in this scenario;
    \item \textbf{Weak Scalability}: Measures how execution time changes when both the processor count $p$ and the total problem size $W$ increase such that the work per processor ($W/p$) stays constant. The aim is to maintain efficiency as $p$ and $W$ grow proportionally. Gustafson’s Law is most relevant here.
\end{itemize}

A parallel system is considered scalable if its efficiency remains above an acceptable threshold as both the number of processors $p$ and the problem size $W$ increase.
% Fim da subseção "Parallel Performance Evaluation Metrics".

% Conteúdo da subseção "Tries (Prefix Trees)"

\subsecao{Tries (Prefix Trees)}

This section explains the Trie data structure, detailing its defining properties, core operational mechanics, performance aspects, and its significance as a foundational element in advanced string processing algorithms, particularly the Aho-Corasick algorithm.

\subsubsecao{Definition and Basic Structure}

A Trie---also known as a prefix tree---is a specialized tree-like data structure optimized for the efficient storage and retrieval of a dynamic set of strings. Edward Fredkin introduced the term ``trie'' in 1960, highlighting its primary application in information retrieval \cite{Fredkin1960}, \cite{Knuth1998}. Figure~\ref{fig:trie-example} illustrates an example of a trie storing a set of words.

% Figura da Trie
% review: comments 009 (issue #104) caption below figure, 010 (issue #105) no source line for author-made figures
\begin{figure}[H]
  \centering % Centraliza a figura
    \includegraphics[width=0.4\textwidth]{trie.pdf}
  \caption{Example trie storing \{``car'', ``cat'', ``cart'', ``dog''\}. Nodes marking the end of a complete word are shown with a double circle.} % Legenda da figura
  \label{fig:trie-example} % Rótulo para referenciar a figura
\end{figure}
% Fim da Figura da Trie

Tries store strings by organizing nodes so that each path from the root to a node represents a unique prefix. In a standard $R$-way trie, each node can branch to $R$ children, where $R$ denotes the size of the character alphabet. The existence of a parent-child edge highlights two key aspects of tries:
\begin{itemize}
    \item The edge's position (index) within the parent node's array of $R$ possible child links directly corresponds to a unique character in the alphabet. For instance, index 0 might correspond to `a', index 1 to `b', and so on, according to a predefined mapping.
    \item The character associated with this edge extends the prefix represented by the parent; this implies that at least one word (or key) in the trie contains this character sequence.
\end{itemize}
To signify that such a prefix also constitutes a complete key present in the trie, the node reached at the end of this path is typically marked by storing an associated non-null value or a boolean flag. The characters themselves are implicitly encoded by the trie's structure rather than being explicitly stored within each node \cite{SedgewickWayne2011}.

\subsubsecao{Core Operations}

Tries support several primary operations, including search, insertion, and deletion \cite{SedgewickWayne2011}. These operations are explained below:

\begin{itemize}
    \item \textbf{Search}: locating a key involves traversing the trie starting from the root, following the successive characters of the target key. Each character dictates which child link to follow. The search operation succeeds if a complete path corresponding to all characters of the key exists and the node reached at the end of this path is appropriately marked as representing a complete key. The search fails if, at any point, a required link is null (i.e., no path for the current character exists) or if the path is fully traversed but the final node is not marked as a complete key;
    \item \textbf{Insertion}: the insertion process begins similarly to a search, by traversing the trie according to the characters of the key to be added. If the path for the key does not extend to its full length (i.e., a null link is encountered), new nodes are created and linked to form the remainder of the path. The node corresponding to the final character of the inserted key is then marked to indicate that it is a complete key in the set;
    \item \textbf{Deletion}: deleting a key first involves locating the node corresponding to the key's final character and removing its end-of-word marker. If there are no child links (meaning it is not a prefix for any other key), it can be deleted. This deletion may proceed recursively: when a child node is removed, if its parent has no other children and the parent node is not marking a key, the parent node becomes redundant and can be deleted. This process continues up to the root as long as parent nodes become redundant.
\end{itemize}

\subsubsecao{Advantages and Performance Characteristics}

Tries provide distinct advantages for operations on string collections, as detailed below.

\begin{itemize}
    \item \textbf{Time Complexity}: a primary advantage is that the time required for search and insertion operations depends primarily on the length of the key ($L$), not on the total number ($N$) of keys stored. Specifically, these operations typically require examining a number of nodes proportional to $L$. Sedgewick and Wayne~\cite{SedgewickWayne2011} state that the number of array accesses (node visits) is at most $L+1$. This $O(L)$ performance makes tries highly efficient for prefix-based searches (e.g., retrieving all keys starting with a given prefix) and ensures that search performance does not substantially degrade as the dataset size increases;
    \item \textbf{Prefix Sharing}: when stored strings share common prefixes, tries inherently exploit this by having those strings share the initial path from the root. This sharing can lead to significant space savings compared to storing each string independently, particularly if there is substantial prefix overlap among keys \cite{Fredkin1960}.
\end{itemize}

While prefix sharing can save space, the overall space complexity of $R$-way tries requires careful consideration. Each node in an $R$-way trie conceptually provides $R$ potential links. If the alphabet size $R$ is large (e.g., for Unicode) and the actual branching factor at most nodes is low (a sparse trie), allocating space for all $R$ links per node can lead to considerable memory consumption due to numerous unused (null) links \cite{Knuth1998}, \cite{SedgewickWayne2011}.

% review: comment 011 (issue #131) — closing wrap-up demoted from a subsection to a reinforcing paragraph
In summary, the trie data structure is a cornerstone for numerous advanced string processing algorithms. Critically for this work, the Aho-Corasick algorithm, engineered for the efficient, simultaneous matching of multiple patterns within a text, employs a trie---constructed from the set of patterns (the dictionary)---as its initial structural framework \cite{AhoCorasick1975}. A thorough grasp of the trie principles presented above is thus essential for developing and analyzing the Aho-Corasick algorithm.

% Fim da subseção "Tries".

% Conteúdo da seção/subseção "Pattern Matching"

\subsecao{Pattern Matching}

This section introduces the fundamental problem of pattern matching within strings, explores key algorithmic approaches to its solution, and highlights its significance across various computational domains. This foundational discussion provides the essential context for understanding advanced multi-pattern matching algorithms, such as the Aho-Corasick algorithm. In general information processing, pattern (or \textit{string}) matching underpins the functionality of text editors and web search engines for keyword searches. The field of computational biology extensively utilizes pattern matching for analyzing genetic sequences \cite{SedgewickWayne2011}. Furthermore, in the field of information security, pattern matching is crucial: Network Intrusion Detection Systems (NIDS) utilize it to identify malicious signatures in network traffic. Tools like YARA, for instance, define rules based on textual or binary patterns to classify and identify malware \cite{YARADocs}.

\subsubsecao{The Pattern Matching Problem}

The \textit{string matching} problem consists of determining whether a string $P$ (the \textit{pattern}) of length $M$ occurs within a string $T$ (the \textit{text}) of length $N$, where $M \leq N$. Formally, it seeks all positions $pos$ such that the condition in Eq.~\ref{eq:string_matching_condition} holds for every $i \in \{0, 1, \dots, M-1\}$, and each valid $pos$ corresponds to an occurrence of the pattern in the text:
\begin{equation}
    P[i] = T[pos + i], \quad \text{with} \quad pos + i < N
    \label{eq:string_matching_condition} 
\end{equation}

Beyond simply finding the first match, the problem often extends to locating all occurrences, enumerating them, or extracting surrounding contextual information for each match \cite{SedgewickWayne2011}.

Algorithms that address the string matching problem are generally categorized based on the number of patterns they process and the exactness of the match. This thesis focuses on \textit{Exact Matching}, where the pattern must correspond precisely to a segment of the text. A key distinction for this study lies between algorithms designed for \textit{Single-Pattern Matching} and those optimized for \textit{Multi-Pattern Matching}, where the objective is to identify all occurrences of any pattern from a predetermined set of patterns within a given text \cite{AhoCorasick1975}, \cite{Gusfield1997}.

\vspace{1em}

\subsubsecao{The Naive Algorithm}

The naive or brute-force algorithm represents the most straightforward method for exact pattern matching. Its operation involves iteratively aligning the pattern $P$ with the text $T$ at every possible starting position $i$ (from 0 to $N-M$). For each alignment, $P$ is compared character by character, from left to right, against the corresponding substring. If all characters match, an occurrence is reported \cite{Gusfield1997}.

Following this comparison (whether a match or a mismatch), the algorithm shifts $P$ one position to the right relative to $T$ and proceeds to the next alignment. Despite its simplicity, Knuth, Morris, and Pratt (1977) state that this method can be highly inefficient \cite{KnuthMorrisPratt1977}. Its worst-case time complexity is $O(M \cdot N)$ character comparisons. This occurs, for example, when searching for a pattern like $P = a^n b$ within a text $T = a^{2n} b$, where many partial matches occur. To address these limitations, more advanced pattern matching algorithms include a preprocessing phase for the pattern $P$. This phase extracts structural information about $P$ that is then utilized during the search, enabling larger shifts of the pattern relative to the text and thus reducing the total number of character comparisons.

\vspace{1em}

\subsubsecao{Knuth-Morris-Pratt (KMP) Algorithm}

The KMP algorithm \cite{KnuthMorrisPratt1977} efficiently resolves the pattern matching problem with an optimal worst-case time complexity of $O(N + M)$. Unlike the na\"ive approach, KMP avoids re-checking characters that have already been matched. To achieve this, KMP first preprocesses the pattern $P$ by constructing a special array known as the \textit{failure array} (or \textit{prefix array}), typically denoted by $p[1..M]$. This array indicates the length of the longest prefix that is also a suffix for each prefix of $P$. With this information, when a mismatch occurs during the search, the algorithm uses the failure array to decide which pattern index to resume comparison from, so the text pointer never backtracks, bypassing unnecessary comparisons and preventing redundant checks.

\vspace{1em}

\subsubsecao{Boyer-Moore (BM) Algorithm}

The Boyer-Moore (BM) algorithm \cite{BoyerMoore1977} is widely recognized for its excellent practical performance, often outperforming linear-time methods in typical cases. Unlike the KMP algorithm, Boyer-Moore typically compares characters starting from the end of the pattern and moves backward.
Its efficiency primarily stems from two rules (or heuristics):
\begin{itemize}
    \item \textbf{Bad Character Rule}: When a mismatch occurs, this rule uses the character from the text that caused the mismatch to determine how far the pattern can safely shift. The algorithm shifts the pattern to align this mismatched text character with its last occurrence in the pattern or moves completely past the mismatched character if it does not appear in the pattern;
    \item \textbf{Good Suffix Rule}: When a mismatch happens after matching one or more characters from the end of the pattern, this rule uses the matched suffix to decide how far to shift the pattern. It searches for another occurrence of the matched suffix within the pattern, ensuring that no previously matched characters are unnecessarily compared again.
\end{itemize}
These two rules combined enable the Boyer-Moore algorithm to skip larger portions of the text, significantly reducing the total number of comparisons.

\subsubsecao{The Need for Specialized Multi-Pattern Algorithms}

When the task involves searching for multiple patterns simultaneously from a set $K = \{P_1, P_2, \dots, P_k\}$ within a text $T$, simply applying a single-pattern matching algorithm for each $P_i$ independently proves highly inefficient. This ``iterative'' approach requires scanning the text multiple times, resulting in a total time complexity that can be approximately $O(k \cdot (N+M_{\text{avg}}))$ or worse, where $M_{\text{avg}}$ is the average pattern length \cite{AhoCorasick1975}.

To efficiently handle applications that require simultaneous searches for a large number of patterns, specialized algorithms are essential. These methods usually preprocess the entire set of patterns to create a unified data structure, which enables a single, efficient pass over the text to identify all occurrences of any pattern within the set \cite{Gusfield1997}.

The theoretical framework of \textit{finite automata} provides a robust and elegant solution to this problem. By constructing a single \textit{finite automaton} (FA) that can recognize all patterns in a given set, the text can be processed in a single continuous scan. Each character read from the text triggers a \textit{state transition} within the FA. When the FA reaches an \textit{accepting state}---a state explicitly designed to indicate the successful recognition of one or more patterns---the algorithm reports an occurrence. Algorithms like the Aho-Corasick algorithm exemplify this approach; they automate the construction of such a \textit{pattern-matching machine}, typically based on a trie structure augmented with specialized \textit{failure transitions} \cite{AhoCorasick1975}. This automaton-based paradigm forms the critical foundation for efficiently solving multi-pattern matching problems. The Aho-Corasick algorithm, which instantiates this paradigm and is the focus of this work, is examined in detail in the next section. % review: comment 013 (issue #133) — forward reference to the next section

% Fim da seção/subseção "Pattern Matching".

% Conteúdo da subseção "The Aho-Corasick Algorithm"

\subsecao{The Aho-Corasick Algorithm}

This section provides a comprehensive examination of the Aho-Corasick algorithm, a highly efficient string-matching method designed to locate all occurrences of multiple patterns (keywords) within a given text. It combines \textit{automata theory} with \textit{trie} data structures to efficiently solve the \textit{multi-pattern search} problem. Its search phase runs in $O(n + z)$ time, where $n$ is the text length and $z$ is the number of matches \cite{AhoCorasick1975}.

\subsubsecao{Context}

The task of concurrently identifying all instances of a large collection of patterns within a text poses a significant computational challenge. A na\"ive approach---running a \textit{single-pattern matching} algorithm such as KMP or BM for each of the $k$ patterns---requires repeated scans of the text and has a worst-case time complexity of $O(k \cdot n)$, where $k$ is the number of patterns and $n$ is the text length, rendering it impractical for large pattern sets or long texts \cite{Gusfield1997}.

The Aho-Corasick algorithm overcomes this by building one \textit{finite automaton} that encodes all patterns simultaneously. This automaton processes the input text in a single pass to detect all matches. The algorithm effectively combines the \textit{state-transition} efficiency of finite automata with the prefix-sharing advantages of tries, following \textit{failure links} on mismatches and resuming matching without rescanning characters---an extension of KMP’s \textit{failure-function} idea to a trie of patterns \cite{Gusfield1997}.

\subsubsecao{Algorithm Components and Construction}

Aho-Corasick has two main phases: (i) construction, which builds the automaton, and (ii) search, which scans the text. The automaton has three key tables, built in the following order:

\begin{itemize}
    \item \textbf{Keyword Trie (goto)}: a \textit{keyword trie} is first constructed from the pattern set $K=\{ P_1, P_2, \dots, P_k \}$. Each node represents a state, and each edge corresponds to a character transition. A path from the root to any state $s$ spells a prefix of one or more keywords. Formally, the \textit{goto function} $\operatorname{goto}(s, a)$ returns the next state $s'$ when an edge labelled $a$ exists; otherwise the transition fails. Undefined transitions from the root are assigned back to the root, i.e., $\operatorname{goto}(0, a) = 0$, creating implicit self-loops for every alphabet symbol;
    \item \textbf{Failure Function (f)}: for any state $s$ (except the root), $\operatorname{f}(s)$ is the state $s'$ whose path label is the longest proper suffix of the path label of $s$ that is also a prefix of some keyword. On a mismatch, the automaton follows $\operatorname{f}(s)$ and resumes matching without consuming the current character. The \textit{failure function} is computed level-by-level (breadth-first) after the goto table is complete;
    \item \textbf{The Output Function}: the \textit{output function}, $\operatorname{output}(s)$, identifies all keywords that end at state $s$. For a given state $s$, $\operatorname{output}(s)$ is the set of all patterns $P_i$ in $K$ such that the string corresponding to the path to $s$ is $P_i$. To ensure all matches are reported, including those that are suffixes of other patterns, $\operatorname{output}(s)$ is augmented during the failure function computation: if $\operatorname{f}(s) = s'$, then $\operatorname{output}(s) \leftarrow \operatorname{output}(s) \cup \operatorname{output}(\operatorname{f}(s))$. This incremental union lets each state report every pattern that ends either at that state or at any suffix state reached through its failure chain.
\end{itemize}

The complete automaton, combining the goto transitions, failure links, and output match sets, is illustrated in Figure~\ref{fig:ac_automaton} for the pattern set $K=\{\text{"he"}, \text{"she"}, \text{"hers"}\}$.

% review: comments 009 (issue #104) caption below figure, 010/017 (issues #105/#135) no source line for author-made figures
\begin{figure}[H]
\centering
\includegraphics[width=1.0\linewidth]{ac_automaton.pdf}
\caption{Complete Aho–Corasick automaton for the pattern set $K=\{\text{"he"}, \text{"she"}, \text{"hers"}\}$, showing solid blue goto transitions, dashed red failure links, double-circled accepting states, and their corresponding output match sets.}
\label{fig:ac_automaton}
\end{figure}


% review: comment 015 (issue #107) — worked example of how the automaton in Figure 3 operates
To make the interaction of the three tables concrete, consider scanning the text \texttt{ushers} with the automaton of Figure~\ref{fig:ac_automaton}. The leading \texttt{u} has no outgoing edge from the root, so the implicit self-loop keeps the machine at state~$0$. The substring \texttt{she} then drives three successive $\operatorname{goto}$ transitions to the accepting state for \texttt{she}; because the failure link of that state points to the state for \texttt{he}---the longest proper suffix of \texttt{she} that is also a keyword prefix---its augmented $\operatorname{output}$ set reports both \texttt{she} and \texttt{he} at that position. The next character \texttt{r} has no $\operatorname{goto}$ edge from the \texttt{she} state, so the automaton follows the failure link to the \texttt{he} state and continues along \texttt{her} and then \texttt{hers}, emitting \texttt{hers} when the final \texttt{s} is consumed. A single left-to-right pass thus reports every occurrence---\texttt{he}, \texttt{she}, and \texttt{hers}---without ever rescanning a character.

\subsubsecao{Matching Phase}

Once the Aho–Corasick automaton has been built, it scans the text $T = a_1 a_2 \dots a_n$ in one left-to-right pass. At each position $i$, the machine attempts to follow the transition labelled $a_i$ from its \textit{current\_state} (cs); if that edge exists, it moves along it, otherwise it follows \textit{failure links} (f) until either a matching transition is found or the root is reached. Every time a successful transition is made, the automaton inspects the $\operatorname{output}$ set of the new state and emits all patterns ending at position $i$.

Figure~\ref{fig:ac-matching-flow} illustrates this control flow in detail, which: starts at the root with $i=1$, attempt each $\operatorname{goto}(\textit{current\_state}, a_i)$, follow failure links on a “no” decision, report matches via $\operatorname{output}(\textit{current\_state})$ on a “yes,” advance $i$, and terminate when $i>n$.

% review: comments 009 (issue #104) caption below figure, 010 (issue #105) no source line for author-made figures
\begin{figure}[H]
    \centering
    \includegraphics[width=0.6\textwidth,keepaspectratio]{ac_matching.pdf}
    \caption{Flowchart of the Aho–Corasick matching phase.}
    \label{fig:ac-matching-flow}
\end{figure}

Each input character advances the automaton by at most one $\operatorname{goto}$ transition, which increases the current depth by one, whereas every failure-link traversal strictly decreases that depth. Since the depth is non-negative, the total number of failure transitions over the scan cannot exceed the number of goto transitions, which bounds the total number of transitions (goto plus failure) by $< 2n$ \cite{AhoCorasick1975}.

\subsubsecao{Time and Space Complexity}

The Aho-Corasick algorithm exhibits optimal \textit{linear time complexity} for multi-pattern matching \cite{AhoCorasick1975}, \cite{Gusfield1997}. The complexities for the construction and search phases and storing the automaton are as follows:

\begin{itemize}
    \item \textbf{Construction phase}: building the goto, failure, and output functions takes $O(L)$ time, where $L$ is the total length of all keywords. If a dense array is used for the root, an additional $O(|\Sigma|)$ initialization cost is incurred, giving $O(L + |\Sigma|)$;
    \item \textbf{Search phase}: Scanning a text of length $n$ performs fewer than $2n$ state transitions. Reporting matches costs time proportional to the total number of occurrences $z$, so the total search time is $O(n + z)$;
    \item \textbf{Space complexity}: Storing the automaton requires $O(L + |\Sigma|)$ space: $O(L)$ for trie edges, $O(L)$ pointers for failure links, and additional space for output lists (proportional to the number of stored outputs).
\end{itemize}

\subsubsecao{NFA versus DFA Representations}

The goto-plus-failure automaton described above is commonly distinguished, in the network-intrusion-detection literature, as the \textit{nondeterministic finite automaton} (NFA) configuration of Aho--Corasick, in contrast to the fully precomputed \textit{deterministic finite automaton} (DFA) discussed below \cite{Qu2016MPI}. This NFA label is conventional rather than formal---the goto-plus-failure machine in fact behaves deterministically---but it captures the operational difference: the goto function is undefined on many (state, character) pairs, and the search algorithm resolves those gaps by traversing failure links at runtime, so a single character may require several memory accesses in the worst case (the cost is constant only in the amortized sense established above). A widely adopted optimization converts this NFA into a \textit{deterministic finite automaton} (DFA) by precomputing the failure-link chains during construction: for every state $s$ and character $a$, the full deterministic transition $\delta(s,a)$ is stored in a flat two-dimensional table of size $|Q| \times |\Sigma|$, where $|Q|$ is the number of states and $|\Sigma|$ is the alphabet size \cite{AhoCorasick1975}, \cite{Gusfield1997}. With this representation every character in the input triggers exactly one array lookup; no runtime failure-link traversal occurs, which reduces per-character latency and simplifies the inner search loop.

The cost of the DFA representation is memory: for the standard 256-symbol byte alphabet the flat transition table occupies $4 \cdot |Q| \cdot 256$ bytes. This footprint grows linearly with the pattern-set size and can easily exceed the processor's last-level cache (LLC). When the table no longer fits in cache, each state transition becomes a potential cache miss, shifting the workload from compute-bound to \textit{memory-bandwidth-bound}. The implementation in this work uses the DFA (flat transition table) representation exclusively; the goto function, failure function, and output function described in the preceding subsection are intermediate construction artefacts that are not retained at search time.

\subsubsecao{Applications and Relevance}

The Aho-Corasick algorithm is widely used in applications that require high-performance \textit{multi-pattern matching}, where its single-pass processing of the input against the entire pattern set is essential. Representative use cases include large-scale text analysis, motif searching in DNA or protein sequences \cite{Vajira2018}, spam filtering, Network Intrusion Detection Systems (NIDS) that scan traffic against thousands of signatures, and anti-malware tooling---YARA, for example, employs Aho-Corasick to match the literal string components extracted from its rules \cite{YARADocs}.

While Aho-Corasick excels at exact literal matching, systems like NIDS often layer additional capabilities---such as \textit{regular-expression matching}---on top of Aho-Corasick engines to support more expressive signatures, which affects overall performance characteristics. The algorithm therefore remains a theoretically elegant and practically impactful example of \textit{finite-automata theory} applied to string processing \cite{Gusfield1997},\cite{YARADocs}.

% Fim da subseção "The Aho-Corasick Algorithm".
